\documentclass[11pt,a4paper]{article}

\usepackage{graphicx}
\usepackage{xcolor}
\usepackage{geometry}
\usepackage{amsmath}
\usepackage{booktabs}
\usepackage{tabularx}
\usepackage{adjustbox}
\usepackage{enumitem}
\usepackage{longtable}
\usepackage{fancyhdr}
\usepackage{titlesec}
\usepackage{listings}
\usepackage{microtype}
\usepackage[colorlinks=true,linkcolor=blue!60!black,citecolor=darkgray,urlcolor=blue!60!black,bookmarks=true,bookmarksnumbered=true,unicode=true]{hyperref}

\geometry{a4paper, top=2.6cm, bottom=2.6cm, left=2.5cm, right=2.5cm}
\setlength{\headheight}{14pt}

% ─── Header / footer ────────────────────────────────────────────────
\pagestyle{fancy}
\fancyhf{}
\fancyhead[L]{\small\itshape covtl-pipeline / vtl-synth --- Wave-4 Commit Log}
\fancyhead[R]{\small base 3903645}
\fancyfoot[C]{\small\thepage}

% ─── Typography ─────────────────────────────────────────────────────
\titleformat{\section}{\Large\bfseries\color{blue!35!black}}{\thesection}{0.8em}{}
\titleformat{\subsection}{\large\bfseries\color{blue!30!black}}{\thesubsection}{0.7em}{}
\setlength{\parskip}{0.45em}
\setlength{\parindent}{0pt}
\widowpenalty=10000
\clubpenalty=10000

\definecolor{codebg}{HTML}{F7F7F5}
\definecolor{okgreen}{HTML}{2E7D32}

\lstset{
  basicstyle=\ttfamily\footnotesize,
  backgroundcolor=\color{codebg},
  frame=single, framerule=0.3pt, rulecolor=\color{gray!40},
  breaklines=true, breakatwhitespace=true,
  columns=fullflexible, keepspaces=true,
  aboveskip=0.8em, belowskip=0.8em,
  showstringspaces=false,
}

\newcommand{\code}[1]{\texttt{#1}}
\newcommand{\commit}[1]{\texttt{#1}}
\newcommand{\pkg}[1]{\texttt{#1}}
\newcommand{\pct}[1]{{#1}\,\%}

\hypersetup{
  pdftitle={Wave-4 Commit Log: Rationale and Non-Regression Evidence (vtl-synth 1.1.3)},
  pdfauthor={Z.ai},
  pdfsubject={Commit-by-commit rationale of the wave-4 patch series with non-regression test series},
  pdfkeywords={git, commits, testing, non-regression, coverage, speech synthesis}
}

\begin{document}

% ════════════════════════════════════════════════════════════════════
% NO TITLE PAGE — cover generated separately, merged as page 0.
% ════════════════════════════════════════════════════════════════════

\tableofcontents
\newpage

% ────────────────────────────────────────────────────────────────────
\section{Purpose and Scope}
\label{sec:purpose}

This document is the engineering record of the wave-4 patch series
applied to \pkg{vtl-synth}~1.1.3 (repository
\pkg{covtl-pipeline}, base commit
\code{3903645d3f91862c49f934aee94e2b2cc09c4bd}). It travels with the
patch: applying the series to the base commit and reading this log
side by side gives a reviewer the intent, the logical justification,
the measured improvement, and the non-regression evidence of every
change, in the order the changes should be reviewed.

The series consists of eight atomic commits, ten files touched,
1{,}170 insertions and 88 deletions. Every commit compiles, passes
the full test suite, and---wherever it is supposed to preserve
behavior---passes the byte-level and structural regression pins of
the repository. The order of the commits is itself part of the design:
the series opens with the test hardening (so that every later commit
is judged by the strengthened suite), proceeds from pure refactoring
(no behavioral change, hash-proven) to bug fixes (intended behavioral
changes, each with a before/after observable), and closes with CI and
coverage-complement commits.

\begin{table}[htbp]
\centering
\caption{The wave-4 series at a glance.}
\label{tab:series}
\small
\begin{tabular*}{\textwidth}{@{\extracolsep{\fill}} c l l l}
\toprule
\# & Commit & Type & Files \\
\midrule
1 & \commit{017acea} & test + ci & 6 (tests, LEGACY, CI, pyproject) \\
2 & \commit{971f2f2} & refactor & 1 (\code{syltraj.py}) \\
3 & \commit{ec81a1a} & fix & 2 (\code{notation.py}, test) \\
4 & \commit{7681cfe} & fix & 2 (\code{enveloppe.py}, test) \\
5 & \commit{269b23a} & fix & 2 (\code{enveloppe.py}, test) \\
6 & \commit{c02f482} & refactor & 3 (\code{enveloppe.py}, \code{build\_phrase\_tract.py}, test) \\
7 & \commit{7d60b6e} & ci & 1 (\code{ci.yml}) \\
8 & \commit{bfc19bb} & test & 1 (\code{test\_wave3\_targeted.py}) \\
\bottomrule
\end{tabular*}
\end{table}

\subsection{How to Apply}

\begin{lstlisting}[language=bash, caption={Applying the series (line endings matter).}]
git checkout -b wave4 3903645
git am --keep-cr covtl-wave4-patch/commits/*.patch
# or, as a single squashed change:
git apply --keep-cr wave4_squashed.diff
\end{lstlisting}

The engine sources of the repository are CRLF blobs in git; the test
sources are LF. The patch series preserves those line endings, and
\code{git am} requires \code{--keep-cr} to replay them (an application
without the flag fails on \code{syltraj.py}---the symptom is a patch
mismatch, not a defect of the series). The series was validated by
applying it to a virgin clone of the base commit: the resulting tree
is byte-identical to the developed branch, and the full suite is
green on the patched tree.

% ────────────────────────────────────────────────────────────────────
\section{Commit 1 --- \commit{017acea}: Wave-3 Test Hardening}
\label{sec:c1}

\subsection{Rationale}

The wave-3 measurement campaign (which preceded this series) had
improved coverage from \pct{74.5} to \pct{76.1} but left its
artifacts---a strengthened phonetic corpus, platform-independent
pins, targeted tests, a CI fast lane, and the coverage
configuration---outside the repository. The first commit of this
series imports that hardening wholesale, so that every subsequent
commit is judged against the strengthened suite rather than the
original one. This is the standard safety pattern for a
behavior-changing series: \emph{raise the bar first}.

\subsection{Content}

\begin{itemize}[leftmargin=1.4em, itemsep=0.15em]
  \item \code{tests/test\_phonetic\_model\_inputs.py}: the phonetic
        corpus grows to 97 test cases, including a batch-load test
        (five real TIMIT phrases) and a label-sanitation uniqueness
        test that locks the silent-dict-overwrite failure mode of
        \code{build\_batch\_tract} (two labels that differ only in
        non-word characters must not collide after sanitation).
  \item \code{tests/test\_syltraj\_double\_pause.py}: the regression
        pins become platform-independent. The single-pause
        trajectories are hashed after quantization to $10^{-9}$
        (removing \code{libm} last-bit differences between Windows
        and Linux while still catching any real trajectory change),
        and structural invariants are added: exact block schedule,
        vowel plateaus equal to the algebraic projection
        \code{compute\_P(*target)} \emph{exactly}, bounded inter-step
        jump. The raw Windows pins are kept as an informational
        control.
  \item \code{tests/test\_wave3\_targeted.py}: 26 targeted tests,
        each citing the exact line it activates (arc variants of
        \code{syltraj}, V1 engine paths, nasal hiatus, speaker-file
        errors, the timed-envelope route,
        \code{\_align\_trajectories} padding and truncation).
  \item \code{LEGACY.md} \S1bis: the hyphen policy---V1 treated
        \code{-} as a syllable separator alias, the syl engine
        silently consumes it---is documented as a locked gap with a
        test that will fail the day someone aligns the behaviors,
        forcing a documentation update.
  \item \code{.github/workflows/ci.yml}: a \code{phonetic-fast} job
        runs the 97-case corpus ($\approx$6\,s) on every push and
        pull request, ahead of the full suite.
  \item \code{pyproject.toml}: the \code{[tool.coverage]} blocks
        (omits with their LEGACY ledger references,
        \code{exclude\_lines}) make local and CI measurements
        identical.
\end{itemize}

\subsection{Evidence}

The commit is test- and infrastructure-only: no engine source is
touched, and the 49 regression baselines pass unchanged by
construction. The corpus run time on this machine: 97 tests in
4.6\,s (wall clock $\approx$5.7\,s including interpreter startup).

% ────────────────────────────────────────────────────────────────────
\section{Commit 2 --- \commit{971f2f2}: Removing the Four Dead Arms}
\label{sec:c2}

\subsection{Rationale}

The wave-2 tracer sweep (\code{sys.settrace} over 25 pause forms) had
established that four control-flow arms of
\code{build\_global\_pval} are structurally unreachable through the
public API; wave~3 had marked them \code{pragma: no cover}. Marking
dead code is an accounting decision; removing it is an engineering
one, and it is the right one: unreachable arms in the central
trajectory builder are a standing trap for future maintainers (every
reader must re-derive their unreachability) and a permanent tax on
the coverage denominator. This commit replaces the marks with
removal---27 statements deleted, each replacement site carrying a
comment that cites the tracer evidence and names the invariant of
\code{build\_anchors} that guarantees unreachability.

\begin{table}[htbp]
\centering
\caption{The four removed arms and the anchor-builder invariant that
makes each unreachable.}
\label{tab:arms}
\small
\begin{tabularx}{\textwidth}{@{ } l X }
\toprule
\textbf{Removed arm} & \textbf{Invariant of \code{build\_anchors}} \\
\midrule
\code{a\_term} \code{else} (cluster after a pause with terminal departure) &
A vowel-onset anchor is inserted after every pause that precedes a consonant, so the consonant count $m$ is always 0 when \code{a\_term} holds; the \code{m==0} arm (arc + \code{initial} block) is the live path. \\
\addlinespace
\code{b\_term} \code{m>0} (terminal cluster) &
A word-end anchor (\code{COEFCEN}$\cdot\rho_V$) is inserted before every pause that follows a consonant, so $m=0$ when \code{b\_term} holds; the \code{elif} (V decay before pause) is live. \\
\addlinespace
Mid-loop pending-pause block &
A non-terminal anchor following a pause is always preceded by a terminal one, so the \code{a\_term} branch consumes the pending point first. \\
\addlinespace
Two post-loop blocks (trailing V plateau, trailing pause arc) &
The anchor chain always terminates with a \code{synth} anchor: never a held vowel, never a leftover pending point. \\
\bottomrule
\end{tabularx}
\end{table}

\subsection{The Error the Harness Caught}

The first attempted edit inverted the live and dead arms of the
\code{a\_term}/\code{b\_term} branches---an easy mistake when the
dead code looks \emph{more} general than the live code. The
bit-identity harness (\code{bitident.py}) compared the fingerprints
of 22 pause phrases before and after the edit and reported 88
differences, including block-schedule changes (\code{['cluster',
...]} where \code{['initial', ...]} was expected). The commit that
was finally recorded is the corrected version, and the incident is
retained here as the operational argument for the harness: a change
of this shape must not be merged on reading alone.

\subsection{Non-Regression Evidence}

\begin{lstlisting}[language=bash, caption={Bit-identity protocol around commit 2.}]
# before the commit (on the parent):
python3 bitident.py save pre
# after applying the commit:
python3 bitident.py check pre
# -> OK - 22 phrases bit-identical (raw + quantized + structure)
\end{lstlisting}

\begin{itemize}[leftmargin=1.4em, itemsep=0.15em]
  \item \textbf{Raw SHA-256}: identical on all 22 pause phrases
        (single and double pauses, breath pauses \code{\%}, TIMIT
        dotted forms, nasal clusters, schwa edge cases).
  \item \textbf{Quantized SHA-256} ($10^{-9}$): identical.
  \item \textbf{Block structure}: kinds, shapes, and step counts of
        every \code{SylBlock} identical.
  \item \textbf{Regression baselines}: the 49 pinned tract/glottis
        SHA-256 tests (schema 3, speaker JD3) pass unchanged.
  \item \textbf{Suite}: all corpus, structure, and double-pause tests
        green (110 tests at commit time).
\end{itemize}

% ────────────────────────────────────────────────────────────────────
\section{Commit 3 --- \commit{ec81a1a}: The Consonant Registry Fix}
\label{sec:c3}

\subsection{Rationale}

\code{notation.ALL\_K} was defined as the union of the vowel and
consonant \emph{target} tables. But the consonant target table
carries only the voiced articulatory targets; the voiceless
counterparts---the \code{\_VOICELESS\_KEYS} of
\code{constants.py}, namely \code{p}, \code{t}, \code{k},
\code{f}, \code{s}, \code{S}, \code{T}, \code{tS}---are
registered in \code{ALL\_PHONEME\_KEYS} and are fully
synthesizable by the engine. Consequence: the reporting helpers
(\code{unknown\_chars}, \code{preprocess\_phrase}) flagged valid,
synthesizable consonants as ``ignored characters'': the phrase
\code{'klinpan'} came back with \code{ignored = ['k', 'p']}
while the tokenizer accepted \code{k} and \code{p} without
complaint---a warning that contradicted the engine.

The fix derives the guard from the complete registry:
\code{ALL\_K = frozenset(ALL\_PHONEME\_KEYS)}. Tokenization is
untouched (the fix only widens the known-character set used for
reporting); what changes is that the report becomes truthful.

\subsection{Measured Improvement}

\begin{itemize}[leftmargin=1.4em, itemsep=0.15em]
  \item Before: \code{'klinpan'} reported ignored
        \code{['k','p']}; the phrase \code{'stra ka pa fjuZ'}
        reported six false positives.
  \item After: both return empty ignored-lists; the accented-vowel
        normalization and the genuinely unknown \code{x} path are
        unchanged (\code{'ba da \'e x'} $\to$ \code{('ba da e x',
        ['x'])}).
  \item \code{core/notation.py} coverage: $\approx$91\,\%
        $\to$ \pct{97.7} (one line missing).
\end{itemize}

\subsection{Non-Regression Evidence}

The 131 notation/corpus/targeted tests pass, including the updated
nominal expectations; the full suite (493 tests at this commit) is
green; the 49 baselines are untouched by construction (the change
affects only the reporting path, not the trajectory path). The
previously documented test that \emph{locked the bug} (expecting
\code{['k','p']}) was rewritten into the nominal expectation---the
commit message records both states so a reviewer sees exactly what
the observable change is.

% ────────────────────────────────────────────────────────────────────
\section{Commit 4 --- \commit{7681cfe}: Repairing \code{parse\_to\_anchors}}
\label{sec:c4}

\subsection{Rationale}

\code{parse\_to\_anchors} is exported in
\code{vtl\_synth.core.\_\_all\_\_}---a public convenience wrapper
combining parsing and anchor construction. It was broken: it
called the parser with the old keyword \code{pause\_marker}
while the parameter had been renamed to \code{coarticulation\_
marker} years of commits earlier; every call raised a
\code{TypeError}. Because no internal code path calls the wrapper,
the API drift had gone unnoticed through every functional test---a
textbook demonstration that exported symbols need their own smoke
tests, not only transitive coverage.

The fix passes the renamed keyword and aligns the defaults with the
callee (300\,ms pause, 60\,ms vowel). The docstring is updated to
name the parameter correctly and to record the drift.

\subsection{Measured Improvement}

\begin{itemize}[leftmargin=1.4em, itemsep=0.15em]
  \item Before: \code{parse\_to\_anchors('ba da')} $\to$
        \code{TypeError}.
  \item After: returns the 8-anchor chain
        \code{[V\_o, C\_b, V, V\_e, V\_o, C\_d, V, V\_e]} with
        strictly non-decreasing anchor times; the single-syllable
        call returns 4 anchors.
  \item Function coverage: \pct{0} (unreachable: raised on entry)
        $\to$ \pct{100}.
\end{itemize}

\subsection{Non-Regression Evidence}

The rewritten nominal test locks the anchor labels, the count, and
the chronological ordering. No engine path uses the wrapper, so the
49 baselines and the bit-identity harness are unaffected by
construction. The 36 targeted and legacy-API tests pass.

% ────────────────────────────────────────────────────────────────────
\section{Commit 5 --- \commit{269b23a}: Tag Overrides Reach the Events}
\label{sec:c5}

\subsection{Rationale}

This fix resolves a defect discovered \emph{during} wave~4,
while trying to cover the phonation-mode event maker: the function
sat at \pct{18.5} with a passing ``breathy'' test whose glottis
delta came from the preset path, not from the feature events.
Root cause: the event scheduler built its tag dictionary from the
default phoneme tags alone, ignoring the level-C overrides carried
by the \code{TagInterpreter} it was given. The override annotated
the segment---the glottal tag was set to \code{'breathy'}---but
the event makers for VOT, voicing, and phonation mode never saw
it. The documented tag-override API (level~C of the three-level
tag system) was silently inoperative for the phonation modes.

The fix mirrors \code{TagInterpreter.apply}: the effective tags are
the defaults updated by the overrides.

\subsection{A Second Layer the Tests Exposed}

With the first fix in place, the coverage still showed the function
body unentered. Investigation found a second mechanism:
\code{TimedEnvelope(anchors=...)} re-derives its segments through
\code{anchors\_to\_segments}, which relabels vowels as
\code{@}---so a level-C override keyed on the original vowel
(\code{'a'}) never matches the converted segment keys. This is not a
bug to fix but a contract to document (the override must target the
converted key), and commit~8 locks it in tests.

\subsection{Measured Improvement}

\begin{itemize}[leftmargin=1.4em, itemsep=0.15em]
  \item Before: a breathy override produced a glottis delta of
        $\approx$0.13 against the modal baseline (i.e.\ nothing); zero
        phonation feature events in the timer.
  \item After: the three modes produce their exact feature-event
        sets---breathy: \code{chink\_area} and \code{x\_bottom};
        creaky: \code{f0\_offset} and \code{double\_pulsing};
        whisper: those plus \code{x\_top}---and the glottis deltas
        are 1.04 / 1.39 / 1.65 respectively.
  \item Event-maker coverage: \pct{18.5} $\to$ \pct{96.3}
        (after commit 8); module \code{enveloppe.py}:
        \pct{81.6} $\to$ \pct{85.3}.
\end{itemize}

\subsection{Non-Regression Evidence}

Default schedules (no overrides) are locked to carry \emph{zero}
phonation events; the timers/expressive suites (90 tests) pass; the
engine outputs are untouched (the fix is on the research-API branch,
not the engine branch), and the baselines remain green.

% ────────────────────────────────────────────────────────────────────
\section{Commit 6 --- \commit{c02f482}: The \code{E\_fallback} Rename}
\label{sec:c6}

\subsection{Rationale}

The two envelope branches of the package overlap on the effort
channel. The branch-B export of \code{TimedEnvelope}
(\code{get\_all\_extensions()}, anchor-indexed) computed an
\code{'E'} trajectory that
\code{build\_phrase\_tract} then \emph{unconditionally overwrote}
with \code{E\_100} (branch~A, block-synchronized, the one true master
of the engine effort) at the junction. The branch-B value was
computed, resampled, and thrown away---dead on arrival, at the cost
of reader confusion about which branch owns the effort.

The fix renames the branch-B export to \code{'E\_fallback'}: an
honest name for a value that is available for research use but is
not what the engine emits. The junction assignment in
\code{build\_phrase\_tract} becomes a creation instead of an
overwrite, with a comment recording the history.

\subsection{Compatibility Statement}

This is the one externally visible rename of the series. Any
third-party consumer reading \code{ext['E']} from
\code{TimedEnvelope.get\_all\_extensions()} must switch to
\code{ext['E\_fallback']}. No internal consumer is affected:
\code{build\_phrase\_tract} and \code{text\_to\_tract} define the
engine \code{'E'} themselves, and \code{assemble\_tract} consumes
only the engine-fed dictionary.

\subsection{Non-Regression Evidence}

The engine outputs are, by construction, unchanged---the overwritten
value never reached them. The 49 baselines and the full suite (493
tests at this commit, including the two updated extension tests that
now assert the presence of \code{'E\_fallback'} and the
\emph{absence} of \code{'E'}) pass.

% ────────────────────────────────────────────────────────────────────
\section{Commit 7 --- \commit{7d60b6e}: The Sandhi CI Lane}
\label{sec:c7}

\subsection{Rationale}

The V1V2 $\times$ sandhi composition test skips whenever only the
Branch-A stub is installed---which is the state of a fresh checkout.
The composition of the two-branch architecture was therefore
\emph{never} verified by continuous integration: it was verified
only in the ad hoc measurement campaigns. The commit adds a dedicated
CI job that closes the gap permanently: clone \pkg{covtl-sandhi},
activate Branch~B with the official copy procedure (copying the
content of \code{sandhi/} over the stub directory
\code{vtl\_synth/utils/sandhi/}---the trailing \code{/.} in the
documented command is what makes it a content copy), assert
\code{SANDHI\_FULL}, then run the crossed suite (hook +
composition).

\subsection{Verification}

The workflow YAML parses (four jobs:
\code{phonetic-fast}, \code{tests}, \code{baselines}, \code{sandhi});
the activation procedure is the one exercised in the wave-3 and
wave-4 measurement campaigns, where the crossed suite runs green
(501 passed, 3 expected skips with Branch~B active).

% ────────────────────────────────────────────────────────────────────
\section{Commit 8 --- \commit{bfc19bb}: Locking the Phonation Modes}
\label{sec:c8}

\subsection{Rationale}

The wave-4 phonation test that accompanied commit~5 drove its
override through the original vowel key (\code{'a'}); because of the
\code{@} relabeling (Section~\ref{sec:c5}), that override never
matched, and the glottis delta the test asserted came from the preset
path. The test \emph{passed} while exercising almost nothing of the
event machinery---the exact failure mode this campaign exists to
eliminate. This commit rewrites it into a contract test: for each of
the three modes, on both a converted vowel key and a consonant key,
the exact feature-event parameter set is asserted; the no-override
baseline is asserted to carry no phonation events at all.

\subsection{Measured Improvement}

\code{\_make\_phonation\_mode\_events}: \pct{18.5} $\to$
\pct{96.3}; the body of the phonation scheduler is now exercised by
a test that says precisely what it expects, rather than by a test
that passes for the wrong reason.

% ────────────────────────────────────────────────────────────────────
\section{Consolidated Non-Regression Test Series}
\label{sec:series}

This section is the checklist a reviewer or release engineer runs
after applying the series. Each command is copy-pasteable; expected
outputs are the ones observed on the developed branch and on the
validated virgin-clone application.

\subsection{Layer 0 --- Patch Integrity (before anything else)}

\begin{lstlisting}[language=bash]
# on a virgin clone of the base commit:
git checkout -b wave4 3903645
git am --keep-cr covtl-wave4-patch/commits/*.patch
git diff wave4-developed 3903645..HEAD   # empty: tree identical
\end{lstlisting}

\subsection{Layer 1 --- Byte-Level Output Pins}

\begin{lstlisting}[language=bash]
pytest tests/test_regression_baselines.py -q
# expected: 49 passed (~3 s)
\end{lstlisting}

\subsection{Layer 2 --- Structural Trajectory Pins (Platform-Independent)}

\begin{lstlisting}[language=bash]
pytest tests/test_syltraj_double_pause.py -q
# expected: all passed; on non-Windows machines an [info] line may
# report that the raw digest differs from the Windows reference pin
# (libm float noise) while the quantized pin matched
\end{lstlisting}

\subsection{Layer 3 --- Bit-Identity Harness (Refactor Proofs)}

\begin{lstlisting}[language=bash]
# with the parent of commit 2 checked out:
python3 bitident.py save pre
git cherry-pick 971f2f2   # or apply patch 0002
python3 bitident.py check pre
# expected: OK - 22 phrases bit-identical (raw + quantized + structure)
\end{lstlisting}

\subsection{Layer 4 --- Functional Contracts (Intended Changes)}

\begin{lstlisting}[language=bash]
pytest tests/test_wave3_targeted.py::test_preprocess_phrase_public_api -q
# consonant fix: 'klinpan' -> ignored []; 'stra ka pa fjuZ' -> ignored []
pytest "tests/test_wave3_targeted.py::test_parse_to_anchors_nominal" -q
# API repair: 8-anchor chain, chronological times
pytest tests/test_wave3_targeted.py -q -k "phonation"
# tag fix: exact event sets for breathy/creaky/whisper
# expected: 26 passed (whole targeted file)
\end{lstlisting}

\subsection{Layer 5 --- Full Suite and Fast Lane}

\begin{lstlisting}[language=bash]
pytest tests/ -q          # 501 passed, 3 skipped (~2-3 min)
pytest tests/test_phonetic_model_inputs.py -q   # 97 passed (~6 s)
\end{lstlisting}

\subsection{Coverage Re-Measurement (Optional Audit)}

\begin{lstlisting}[language=bash]
pytest tests/ --cov=vtl_synth --cov-branch \
  --cov-report=term-missing:skip-covered \
  --cov-report=xml:coverage_wave4.xml -q
# constant scope (sandhi excluded): 80.2% lines, 67.0% branches,
# 76.6% combined; honest ceiling 78.6% combined
\end{lstlisting}

\section{Closing Notes for Reviewers}

The series was designed so that each commit can be reverted
independently if a maintainer disagrees with a decision: the dead-arm
removal stands alone (hash-proven), the three fixes are orthogonal
to each other, and the CI/test commits touch no engine source. The
one commit that requires coordinated action is the
\code{'E'}~$\to$~\code{'E\_fallback'} rename if downstream consumers
exist---hence its compatibility statement in Section~\ref{sec:c6}.

Every number quoted in this log was re-verified against the
campaign's coverage XML at document-writing time; the XML ships
with the patch. The companion synthesis report (same series,
directory \code{docs/wave4/synthesis/}, file
\code{COVERAGE\_SYNTHESIS\_WAVE1-4.pdf}) carries the full
four-wave tables and the PyPI promotion case.

\end{document}
